\documentclass[10pt,a4paper]{amsart}
\usepackage{amsmath,amssymb,amsthm,graphicx}
\usepackage{verbatim}
\newenvironment{reference}{\comment}{\endcomment}
%\newenvironment{reference}{}{}
\newenvironment{slogan}{\comment}{\endcomment}
\newenvironment{history}{\comment}{\endcomment}
% For commutative diagrams we use Xy-pic
\usepackage[all]{xy}

% We use 2cell for 2-commutative diagrams.
\xyoption{2cell}
% Theorem environments.
%
\theoremstyle{plain}
\newtheorem{theorem}[subsection]{Theorem}
\newtheorem{proposition}[subsection]{Proposition}
\newtheorem{lemma}[subsection]{Lemma}

\theoremstyle{definition}
\newtheorem{definition}[subsection]{Definition}
\newtheorem{example}[subsection]{Example}
\newtheorem{exercise}[subsection]{Exercise}
\newtheorem{situation}[subsection]{Situation}

\theoremstyle{remark}
\newtheorem{remark}[subsection]{Remark}
\newtheorem{remarks}[subsection]{Remarks}

\numberwithin{equation}{subsection}

% Macros
%
\def\lim{\mathop{\mathrm{lim}}\nolimits}
\def\colim{\mathop{\mathrm{colim}}\nolimits}
\def\Spec{\mathop{\mathrm{Spec}}}
\def\Hom{\mathop{\mathrm{Hom}}\nolimits}
\def\Ext{\mathop{\mathrm{Ext}}\nolimits}
\def\SheafHom{\mathop{\mathcal{H}\!\mathit{om}}\nolimits}
\def\SheafExt{\mathop{\mathcal{E}\!\mathit{xt}}\nolimits}
\def\Sch{\mathit{Sch}}
\def\Mor{\mathop{\mathrm{Mor}}\nolimits}
\def\Ob{\mathop{\mathrm{Ob}}\nolimits}
\def\Sh{\mathop{\mathit{Sh}}\nolimits}
\def\NL{\mathop{N\!L}\nolimits}
\def\CH{\mathop{\mathrm{CH}}\nolimits}
\def\proetale{{pro\text{-}\acute{e}tale}}
\def\etale{{\acute{e}tale}}
\def\QCoh{\mathit{QCoh}}
\def\Ker{\mathop{\mathrm{Ker}}}
\def\Im{\mathop{\mathrm{Im}}}
\def\Coker{\mathop{\mathrm{Coker}}}
\def\Coim{\mathop{\mathrm{Coim}}}

% Boxtimes
%
\DeclareMathSymbol{\boxtimes}{\mathbin}{AMSa}{"02}

%
% Macros for moduli stacks/spaces
%
\def\QCohstack{\mathcal{QC}\!\mathit{oh}}
\def\Cohstack{\mathcal{C}\!\mathit{oh}}
\def\Spacesstack{\mathcal{S}\!\mathit{paces}}
\def\Quotfunctor{\mathrm{Quot}}
\def\Hilbfunctor{\mathrm{Hilb}}
\def\Curvesstack{\mathcal{C}\!\mathit{urves}}
\def\Polarizedstack{\mathcal{P}\!\mathit{olarized}}
\def\Complexesstack{\mathcal{C}\!\mathit{omplexes}}
% \Pic is the operator that assigns to X its picard group, usage \Pic(X)
% \Picardstack_{X/B} denotes the Picard stack of X over B
% \Picardfunctor_{X/B} denotes the Picard functor of X over B
\def\Pic{\mathop{\mathrm{Pic}}\nolimits}
\def\Picardstack{\mathcal{P}\!\mathit{ic}}
\def\Picardfunctor{\mathrm{Pic}}
\def\Deformationcategory{\mathcal{D}\!\mathit{ef}}
\usepackage{cleveref}
\crefname{originallemma}{original lemma}{original lemmas}
\Crefname{originallemma}{Original lemma}{Original lemmas}
\makeatletter
\newlabel{lemma-tensor-product-permanence}{{lemma-tensor-product-permanence}{source}{}{sourcepointer-1}{}}
\newlabel{lemma-tensor-product-permanence@cref}{{[originallemma][0][]lemma-tensor-product-permanence}{[1][0][]source}{}{}{}}
\makeatother
\crefname{originallemma}{original lemma}{original lemmas}
\Crefname{originallemma}{Original lemma}{Original lemmas}
\makeatletter
\newlabel{lemma-presentation-sym-exterior}{{lemma-presentation-sym-exterior}{source}{}{sourcepointer-2}{}}
\newlabel{lemma-presentation-sym-exterior@cref}{{[originallemma][0][]lemma-presentation-sym-exterior}{[1][0][]source}{}{}{}}
\makeatother
\crefname{originallemma}{original lemma}{original lemmas}
\Crefname{originallemma}{Original lemma}{Original lemmas}
\makeatletter
\newlabel{lemma-local-tensor-algebra}{{lemma-local-tensor-algebra}{source}{}{sourcepointer-3}{}}
\newlabel{lemma-local-tensor-algebra@cref}{{[originallemma][0][]lemma-local-tensor-algebra}{[1][0][]source}{}{}{}}
\makeatother
\crefname{originallemma}{original lemma}{original lemmas}
\Crefname{originallemma}{Original lemma}{Original lemmas}
\makeatletter
\newlabel{lemma-coherent-abelian}{{lemma-coherent-abelian}{source}{}{sourcepointer-4}{}}
\newlabel{lemma-coherent-abelian@cref}{{[originallemma][0][]lemma-coherent-abelian}{[1][0][]source}{}{}{}}
\makeatother
\crefname{originalreference}{original reference}{original references}
\Crefname{originalreference}{Original reference}{Original references}
\makeatletter
\newlabel{algebra-lemma-free-tensor-algebra}{{algebra-lemma-free-tensor-algebra}{source}{}{sourcepointer-5}{}}
\newlabel{algebra-lemma-free-tensor-algebra@cref}{{[originalreference][0][]algebra-lemma-free-tensor-algebra}{[1][0][]source}{}{}{}}
\makeatother
\crefname{originalreference}{original reference}{original references}
\Crefname{originalreference}{Original reference}{Original references}
\makeatletter
\newlabel{algebra-lemma-present-sym-wedge}{{algebra-lemma-present-sym-wedge}{source}{}{sourcepointer-6}{}}
\newlabel{algebra-lemma-present-sym-wedge@cref}{{[originalreference][0][]algebra-lemma-present-sym-wedge}{[1][0][]source}{}{}{}}
\makeatother
\title{Local generation, finite type, finite presentation, coherence, quasi-coherence and local freeness pass to tensor powers, exterior powers and symmetric powers}
\author{The Stacks Project Authors}
\date{stacks-project:modules.tex:local-snapshot-20261012; GNU FDL 1.2 or later}
\begin{document}
\maketitle
\noindent\textit{Editorial note.} All delivered mathematical text is the author's own, copied from the pinned Stacks Project chapter source sources/stacks-modules/pinned-original-source.tex; no statement or proof step was rewritten, reordered or removed.\par\medskip
\noindent\textit{Editorial note.} Preamble provenance: the Stacks Project keeps one single shared preamble for every chapter; this card takes it directly from the staged sources/stacks-modules/preamble.tex as three verbatim spans (lines 8-12, 14-18 and 154-217), so no derived file was written and nothing was added or redefined. The wrapper class and the base packages are supplied by the assembler.\par\medskip
\noindent\textit{Editorial note.} Bibliography provenance: neither the counted statement nor its proof contains a citation (this lane rescans the whole span before writing the card), so the card carries no bibliography and no entry was reconstructed.\par\medskip
\noindent\textit{Editorial note.} Selection provenance: this card was opened under the file-based control contract (CONTROL.md), with the mechanical pre-card gate run on its statement and proof interval before the card was written.\par\medskip
\noindent\textit{Allowed standard prerequisites (editorial).}\begin{itemize}
\item The corresponding permanence statements for the tensor powers themselves. Invoked at the start of the counted proof (source lines 3627-3628, 3646-3647) as a pointer to a lemma of this chapter. Reduces all six assertions to statements about the tensor powers, the exterior and symmetric powers being quotients of them.
\item The presentation of the symmetric and exterior powers and of their tensor powers by generators and relations, and the behaviour of coherence under exact sequences of modules. Invoked inside the counted proof (source lines 3640-3647, 3650-3657, 3669-3673) as pointers to lemmas of this chapter and to results of another chapter. Produces the presentation used to transfer finite presentation and coherence.
\end{itemize}\medskip
\ifdefined\hypertarget\hypertarget{sourcepointer-1}{}\fi
\noindent\textit{Source reference.} Original lemma lemma-tensor-product-permanence: sources/stacks-modules/pinned-original-source.tex, source lines 3628--3628. Navigation pointer only: the counted text refers to the Stacks item lemma-tensor-product-permanence, which belongs to another chapter of the project and is not part of this card. The pointer carries the target's own label as its number because this chapter's staged snapshot does not contain the other chapter's tag table; its content is not reproduced here.\par\medskip
\ifdefined\hypertarget\hypertarget{sourcepointer-2}{}\fi
\noindent\textit{Source reference.} Original lemma lemma-presentation-sym-exterior: sources/stacks-modules/pinned-original-source.tex, source lines 3641--3641. Navigation pointer only: the counted text refers to the Stacks item lemma-presentation-sym-exterior, which belongs to another chapter of the project and is not part of this card. The pointer carries the target's own label as its number because this chapter's staged snapshot does not contain the other chapter's tag table; its content is not reproduced here.\par\medskip
\ifdefined\hypertarget\hypertarget{sourcepointer-3}{}\fi
\noindent\textit{Source reference.} Original lemma lemma-local-tensor-algebra: sources/stacks-modules/pinned-original-source.tex, source lines 3551--3551. Navigation pointer only: the counted text refers to the Stacks item lemma-local-tensor-algebra, which belongs to another chapter of the project and is not part of this card. The pointer carries the target's own label as its number because this chapter's staged snapshot does not contain the other chapter's tag table; its content is not reproduced here.\par\medskip
\ifdefined\hypertarget\hypertarget{sourcepointer-4}{}\fi
\noindent\textit{Source reference.} Original lemma lemma-coherent-abelian: sources/stacks-modules/pinned-original-source.tex, source lines 2011--2011. Navigation pointer only: the counted text refers to the Stacks item lemma-coherent-abelian, which belongs to another chapter of the project and is not part of this card. The pointer carries the target's own label as its number because this chapter's staged snapshot does not contain the other chapter's tag table; its content is not reproduced here.\par\medskip
\ifdefined\hypertarget\hypertarget{sourcepointer-5}{}\fi
\noindent\textit{Source reference.} Original reference algebra-lemma-free-tensor-algebra: sources/stacks-modules/pinned-original-source.tex, source lines 3637--3637. Navigation pointer only: the counted text refers to the Stacks item algebra-lemma-free-tensor-algebra, which belongs to another chapter of the project and is not part of this card. The pointer carries the target's own label as its number because this chapter's staged snapshot does not contain the other chapter's tag table; its content is not reproduced here.\par\medskip
\ifdefined\hypertarget\hypertarget{sourcepointer-6}{}\fi
\noindent\textit{Source reference.} Original reference algebra-lemma-present-sym-wedge: sources/stacks-modules/pinned-original-source.tex, source lines 3651--3651. Navigation pointer only: the counted text refers to the Stacks item algebra-lemma-present-sym-wedge, which belongs to another chapter of the project and is not part of this card. The pointer carries the target's own label as its number because this chapter's staged snapshot does not contain the other chapter's tag table; its content is not reproduced here.\par\medskip
\begin{lemma}
\label{lemma-tensor-algebra-permanence}
Let $(X, \mathcal{O}_X)$ be a ringed space.
Let $\mathcal{F}$ be a sheaf of $\mathcal{O}_X$-modules.
\begin{enumerate}
\item If $\mathcal{F}$ is locally generated by sections,
then so is each $\text{T}^n(\mathcal{F})$,
$\wedge^n(\mathcal{F})$, and $\text{Sym}^n(\mathcal{F})$.
\item If $\mathcal{F}$ is of finite type,
then so is each $\text{T}^n(\mathcal{F})$,
$\wedge^n(\mathcal{F})$, and $\text{Sym}^n(\mathcal{F})$.
\item If $\mathcal{F}$ is of finite presentation,
then so is each $\text{T}^n(\mathcal{F})$,
$\wedge^n(\mathcal{F})$, and $\text{Sym}^n(\mathcal{F})$.
\item If $\mathcal{F}$ is coherent,
then for $n > 0$ each $\text{T}^n(\mathcal{F})$,
$\wedge^n(\mathcal{F})$, and $\text{Sym}^n(\mathcal{F})$
is coherent.
\item If $\mathcal{F}$ is quasi-coherent,
then so is each $\text{T}^n(\mathcal{F})$,
$\wedge^n(\mathcal{F})$, and $\text{Sym}^n(\mathcal{F})$.
\item If $\mathcal{F}$ is locally free,
then so is each $\text{T}^n(\mathcal{F})$,
$\wedge^n(\mathcal{F})$, and $\text{Sym}^n(\mathcal{F})$.
\end{enumerate}
\end{lemma}

\begin{proof}
These statements for $\text{T}^n(\mathcal{F})$ follow
from Lemma \ref{lemma-tensor-product-permanence}.

\medskip\noindent
Statements (1) and (2) follow from the fact that
$\wedge^n(\mathcal{F})$ and $\text{Sym}^n(\mathcal{F})$
are quotients of $\text{T}^n(\mathcal{F})$.

\medskip\noindent
Statement (6) follows from
Algebra, Lemma \ref{algebra-lemma-free-tensor-algebra}.

\medskip\noindent
For (3) and (5) we will use
Lemma \ref{lemma-presentation-sym-exterior} above.
By locally choosing a presentation
$\mathcal{F}_2 \to \mathcal{F}_1 \to \mathcal{F} \to 0$
with $\mathcal{F}_i$ free, or finite free and applying the
lemma we see that $\text{Sym}^n(\mathcal{F})$, $\wedge^n(\mathcal{F})$
has a similar presentation; here we use (6) and
Lemma \ref{lemma-tensor-product-permanence}.

\medskip\noindent
To prove (4) we will use
Algebra, Lemma \ref{algebra-lemma-present-sym-wedge}.
We may localize on $X$ and assume that
$\mathcal{F}$ is generated by a finite set
$(s_i)_{i \in I}$ of global sections.
The lemma mentioned above
combined with Lemma \ref{lemma-local-tensor-algebra} above
implies that for $n \geq 2$
there exists an exact sequence
$$
\bigoplus\nolimits_{j \in J}
\text{T}^{n - 2}(\mathcal{F})
\to
\text{T}^n(\mathcal{F})
\to
\text{Sym}^n(\mathcal{F})
\to
0
$$
where the index set $J$ is finite. Now we know that
$\text{T}^{n - 2}(\mathcal{F})$ is finitely generated
and hence the image of the first arrow is a coherent
subsheaf of $\text{T}^n(\mathcal{F})$, see Lemma \ref{lemma-coherent-abelian}.
By that same lemma we conclude that $\text{Sym}^n(\mathcal{F})$ is
coherent.
\end{proof}

\end{document}
